% ============================================================
% Звіт до проєкту
% Дисципліна: Інноваційні сфери застосування нейронних мереж
% КПІ ім. Ігоря Сікорського — ІПСА — Кафедра ШІ
% ============================================================
\documentclass[12pt,a4paper]{article}

\usepackage[T2A]{fontenc}
\usepackage[utf8]{inputenc}
\usepackage[ukrainian]{babel}
\usepackage[a4paper, top=2cm, bottom=2cm, left=2.5cm, right=1.5cm]{geometry}
\usepackage{booktabs}
\usepackage{multirow}
\usepackage{rotating}
\usepackage{array}
\usepackage{amsmath}
\usepackage{setspace}
\usepackage{indentfirst}
\usepackage{caption}
\usepackage[colorlinks=true, linkcolor=black, citecolor=black, urlcolor=blue]{hyperref}
\usepackage{microtype}

\hyphenpenalty=10000
\exhyphenpenalty=10000

\setlength{\parindent}{1.25cm}
\onehalfspacing
\setlength{\emergencystretch}{3em}
\pagestyle{plain}
\captionsetup{font=small, labelfont=bf}

% ---- helpers for tables: best = bold, second = underline ----
\newcommand{\best}[1]{\textbf{#1}}
\newcommand{\second}[1]{\underline{#1}}

\begin{document}

% ============================================================
% TITLE PAGE
% ============================================================
\begin{titlepage}
\begin{center}
\textbf{НАЦІОНАЛЬНИЙ ТЕХНІЧНИЙ УНІВЕРСИТЕТ УКРАЇНИ}\\[0.2cm]
\textbf{«КИЇВСЬКИЙ ПОЛІТЕХНІЧНИЙ ІНСТИТУТ}\\
\textbf{ІМЕНІ ІГОРЯ СІКОРСЬКОГО»}\\[0.5cm]
Інститут прикладного системного аналізу\\
Кафедра штучного інтелекту\\[2.5cm]

{\large\bfseries ЗВІТ ДО ПРОЄКТУ}\\[0.5cm]

{\large з дисципліни «Інноваційні сфери застосування нейронних мереж»}\\[0.8cm]

{\large\bfseries на тему:}\\[0.3cm]
{\large «Аналіз нейромережевих послідовних та базових статичних\\
методів у задачах колаборативної фільтрації»}\\[3.5cm]
\end{center}

\begin{flushright}
\begin{tabular}{rl}
Виконав:      & аспірант групи [ГРУПА] \\
              & [ПІБ] \\[1.2em]
Перевірила:   & [ПІБ] \\
\end{tabular}
\end{flushright}

\vfill
\begin{center}
Київ, 2026
\end{center}
\end{titlepage}

% ============================================================
% SECTION 1 — МЕТА РОБОТИ
% ============================================================
\section{Мета роботи}

Метою даного проєкту є систематичне порівняння восьми методів
колаборативної фільтрації (послідовних та загальних) на єдиній
відтворюваній методологічній основі.
Конкретні завдання:
\begin{itemize}
    \item оцінити якість рекомендацій кожної моделі за повним набором метрик
          (HR, NDCG, MRR, Precision) при порогах $K\!\in\!\{10,20,50,100\}$;
    \item кількісно визначити розрив між послідовними та загальними підходами;
    \item встановити оптимальну модель за критерієм «якість\,/\,час навчання».
\end{itemize}

Набір даних: MovieLens-1M~\cite{harper2015movielens}.
Реалізації моделей взяті з фреймворку RecBole~\cite{zhao2021recbole};
гіперпараметри підбираються автоматично за допомогою Optuna~\cite{akiba2019optuna}.

% ============================================================
% SECTION 2 — НАБІР ДАНИХ
% ============================================================
\section{Набір даних}

MovieLens-1M~\cite{harper2015movielens} є загальновживаним еталонним набором
даних платформи GroupLens Research, що містить оцінки фільмів, зібрані з
кінця 1990-х по 2003~рік.
У даному дослідженні оцінки трактуються як сигнал неявного зворотного
зв'язку (факт перегляду).
Застосовано фільтрацію: залишаються лише користувачі та фільми з
щонайменше~5 взаємодіями.
Характеристики після фільтрації наведено в таблиці~\ref{tab:dataset}.

\begin{table}[h!]
\centering
\caption{Характеристики набору даних MovieLens-1M після фільтрації.}
\label{tab:dataset}
\begin{tabular}{lr}
\toprule
Параметр & Значення \\
\midrule
Кількість взаємодій   & 1\,000\,209 \\
Кількість користувачів & 6\,040 \\
Кількість фільмів      & 3\,706 \\
Щільність матриці      & 4{,}47\% \\
\bottomrule
\end{tabular}
\end{table}

% ============================================================
% SECTION 3 — МЕТОДОЛОГІЯ
% ============================================================
\section{Методологія}

\subsection{Порівнювані моделі}

Досліджуються вісім моделей двох сімейств (таблиця~\ref{tab:models}).
\textit{Загальні} методи (Pop, BPR~\cite{rendle2009bpr},
ItemKNN~\cite{sarwar2001itemknn}) не враховують порядок взаємодій і
є базовими моделями.
\textit{Послідовні} методи явно моделюють хронологічну послідовність:
FPMC~\cite{rendle2010fpmc} поєднує матричну факторизацію з марківськими
ланцюгами;
GRU4Rec~\cite{hidasi2016gru4rec} використовує рекурентну мережу GRU;
NARM~\cite{li2017narm} доповнює GRU механізмом уваги;
SASRec~\cite{kang2018sasrec} побудований на авторегресивній трансформерній
увазі;
BERT4Rec~\cite{sun2019bert4rec} адаптує двонаправлену увагу з маскуванням токенів
BERT для задачі рекомендацій.

\begin{table}[h!]
\centering
\caption{Порівнювані моделі та їх ключові характеристики.}
\label{tab:models}
\begin{tabular}{llp{7.5cm}}
\toprule
Модель   & Сімейство  & Ключова архітектурна ідея \\
\midrule
Pop      & загальна   & ранжування за глобальною популярністю \\
BPR      & загальна   & парне ранжування з імпліцитним сигналом \\
ItemKNN  & загальна   & item-based колаборативна фільтрація \\
\midrule
FPMC     & послідовна & матрична факторизація + ланцюги Маркова \\
GRU4Rec  & послідовна & рекурентна мережа (GRU) \\
NARM     & послідовна & GRU + механізм м'якої уваги \\
SASRec   & послідовна & авторегресивна трансформерна увага \\
BERT4Rec & послідовна & двонапрямлена увага з маскуванням токенів \\
\bottomrule
\end{tabular}
\end{table}

\subsection{Розбиття даних}

Для кожного користувача взаємодії впорядковуються за часовою міткою.
Остання взаємодія відводиться до тестової вибірки, передостання
до валідаційної, а всі попередні до навчальної
(\textit{leave-one-out} стратегія; параметри RecBole:
\texttt{LS=valid\_and\_test}, \texttt{order=TO}).
Завданням моделі є правильно передбачити наступний фільм у персональній
хронологічній послідовності користувача (next-item prediction).

\subsection{Підбір гіперпараметрів}

Автоматичний підбір гіперпараметрів здійснюється бібліотекою
Optuna~\cite{akiba2019optuna} з алгоритмом пошуку TPE та MedianPruner.
Для кожної моделі виконується 10~пробних запусків по 10~епох
(HPO\_TRIALS\,=\,10, HPO\_EPOCHS\,=\,10); фінальна модель навчається
50~епох (FINAL\_EPOCHS\,=\,50) з найкращою комбінацією гіперпараметрів.
Стани досліджень зберігаються у SQLite-базі, що забезпечує можливість
відновлення при перериванні.

\subsection{Метрики оцінювання}

Оцінювання проводиться у режимі повного словника
(\textit{full-vocabulary scoring}): кожна модель ранжує всі фільми набору
даних для кожного тестового користувача, попередньо маскуючи об'єкти з
його навчальної та валідаційної історії.
Обчислюються чотири групи метрик при порогах $K\!\in\!\{10,20,50,100\}$:

\begin{itemize}
    \item \textbf{HR@K} (Hit Rate): частка користувачів, для яких цільовий
          фільм потрапив до топ-$K$;
    \item \textbf{NDCG@K}: нормалізований дисконтований кумулятивний виграш,
          враховує позицію релевантного результату в ранжованому списку;
    \item \textbf{MRR@K}: середній зворотній ранг першого релевантного
          результату;
    \item \textbf{Precision@K}: частка релевантних об'єктів серед топ-$K$.
\end{itemize}

\noindent Оскільки у leave-one-out оцінюванні кожен користувач має рівно один
тестовий об'єкт, $\text{HR@K}\equiv\text{Recall@K}$, тому Recall окремо
не наводиться.

% ============================================================
% SECTION 4 — РЕЗУЛЬТАТИ
% ============================================================
\section{Результати}

Повна таблиця результатів наведена в таблиці~\ref{tab:results}
(горизонтальна орієнтація).
Порівняння середніх показників двох сімейств наведено в
таблиці~\ref{tab:family}.

% --- WIDE LANDSCAPE TABLE ---
\begin{sidewaystable}[p]
\centering
\scriptsize
\setlength{\tabcolsep}{3.8pt}
\caption{Результати оцінювання на тестовій вибірці MovieLens-1M.
         HR\,=\,Recall за протоколом leave-one-out.
         \best{Жирним} виділено найкраще значення у стовпці,
         \underline{підкресленням} позначено друге.
         Train: час навчання фінальної моделі, с.}
\label{tab:results}
\begin{tabular}{l|rrrr|rrrr|rrrr|rrrr|r}
\toprule
& \multicolumn{4}{c|}{\textbf{HR@K}}
& \multicolumn{4}{c|}{\textbf{NDCG@K}}
& \multicolumn{4}{c|}{\textbf{MRR@K}}
& \multicolumn{4}{c|}{\textbf{Precision@K}}
& \\
\textbf{Модель}
  & \multicolumn{1}{c}{10}
  & \multicolumn{1}{c}{20}
  & \multicolumn{1}{c}{50}
  & \multicolumn{1}{c|}{100}
  & \multicolumn{1}{c}{10}
  & \multicolumn{1}{c}{20}
  & \multicolumn{1}{c}{50}
  & \multicolumn{1}{c|}{100}
  & \multicolumn{1}{c}{10}
  & \multicolumn{1}{c}{20}
  & \multicolumn{1}{c}{50}
  & \multicolumn{1}{c|}{100}
  & \multicolumn{1}{c}{10}
  & \multicolumn{1}{c}{20}
  & \multicolumn{1}{c}{50}
  & \multicolumn{1}{c|}{100}
  & \textbf{Train, с} \\
\midrule
Pop
  & 0.0368 & 0.0687 & 0.1387 & 0.2257
  & 0.0177 & 0.0257 & 0.0395 & 0.0535
  & 0.0120 & 0.0142 & 0.0164 & 0.0176
  & 0.0037 & 0.0034 & 0.0028 & 0.0023
  & 8.5 \\
BPR
  & 0.0684 & 0.1210 & 0.2278 & 0.3609
  & 0.0341 & 0.0474 & 0.0685 & 0.0901
  & 0.0239 & 0.0275 & 0.0309 & 0.0328
  & 0.0068 & 0.0061 & 0.0046 & 0.0036
  & 91.3 \\
ItemKNN
  & 0.0745 & 0.1268 & 0.2422 & 0.3639
  & 0.0372 & 0.0503 & 0.0730 & 0.0927
  & 0.0261 & 0.0296 & 0.0332 & 0.0349
  & 0.0075 & 0.0063 & 0.0048 & 0.0036
  & 21.3 \\
\midrule
FPMC
  & 0.1599 & 0.2349 & 0.3738 & 0.5063
  & 0.0813 & 0.1003 & 0.1279 & 0.1494
  & 0.0574 & 0.0626 & 0.0671 & 0.0690
  & 0.0160 & 0.0117 & 0.0075 & 0.0051
  & 514.6 \\
GRU4Rec
  & \best{0.2467} & \best{0.3455} & 0.5017 & 0.6151
  & \best{0.1436} & \best{0.1685} & \best{0.1994} & \best{0.2178}
  & \best{0.1121} & \best{0.1188} & \best{0.1238} & \best{0.1254}
  & \best{0.0247} & \best{0.0173} & 0.0100 & \second{0.0062}
  & 644.7 \\
NARM
  & 0.2247 & 0.3253 & 0.4891 & 0.6127
  & 0.1163 & 0.1416 & 0.1741 & 0.1942
  & 0.0834 & 0.0903 & 0.0955 & 0.0973
  & 0.0225 & 0.0163 & 0.0098 & 0.0061
  & 468.3 \\
SASRec
  & \second{0.2444} & 0.3430 & \best{0.5119} & \best{0.6313}
  & \second{0.1272} & \second{0.1521} & \second{0.1857} & \second{0.2051}
  & \second{0.0916} & \second{0.0984} & \second{0.1038} & \second{0.1055}
  & \second{0.0244} & \second{0.0172} & \best{0.0102} & \best{0.0063}
  & 1952.1 \\
BERT4Rec
  & 0.2359 & \second{0.3454} & \second{0.5063} & \second{0.6248}
  & 0.1237 & 0.1512 & 0.1832 & 0.2024
  & 0.0898 & 0.0973 & 0.1024 & 0.1041
  & 0.0236 & \best{0.0173} & \second{0.0101} & \second{0.0062}
  & 2473.0 \\
\bottomrule
\end{tabular}
\end{sidewaystable}

\begin{table}[h!]
\centering
\caption{Порівняння двох сімейств моделей (середні значення по групі).}
\label{tab:family}
\begin{tabular}{lrrr}
\toprule
Сімейство & NDCG@10 & HR@10 & MRR@10 \\
\midrule
Загальні (Pop, BPR, ItemKNN) & 0.0297 & 0.0599 & 0.0207 \\
Послідовні (FPMC, GRU4Rec, NARM, SASRec, BERT4Rec) & 0.1184 & 0.2223 & 0.0869 \\
\midrule
Відношення (послідовні / загальні) & $\times$3.99 & $\times$3.71 & $\times$4.20 \\
\bottomrule
\end{tabular}
\end{table}

% ============================================================
% SECTION 5 — АНАЛІЗ РЕЗУЛЬТАТІВ
% ============================================================
\section{Аналіз результатів}

\paragraph{Перевага послідовних моделей.}
Таблиця~\ref{tab:family} демонструє різку різницю між двома сімействами: в
середньому послідовні моделі перевищують загальні у 3.7--4.2~рази за
ключовими метриками.
Це підтверджує, що хронологічна структура взаємодій несе суттєву
прогностичну інформацію, яку статичні методи не здатні використати.
Навіть найпростіша послідовна модель (FPMC, NDCG@10\,=\,0.0813) вдвічі
перевершує найкращу загальну базову модель (ItemKNN, NDCG@10\,=\,0.0372).

\paragraph{Лідер: GRU4Rec.}
Рекурентна модель GRU4Rec досягає найвищих значень за NDCG@10\,=\,0.1436,
HR@10\,=\,0.2467, MRR@10\,=\,0.1121 та всіма метриками при $K\le 50$.
При цьому час навчання становить лише 644~с, що є помірним: майже утричі
менший, ніж у SASRec (1952~с), і майже вчетверо менший, ніж у BERT4Rec (2473~с).
Такий результат свідчить про те, що GRU ефективно вловлює короткострокові
залежності у послідовностях взаємодій без надмірних витрат на навчання.

\paragraph{Трансформерні моделі: SASRec та BERT4Rec.}
SASRec та BERT4Rec є конкурентоспроможними (NDCG@10 відповідно 0.1272 та
0.1237) і перевершують GRU4Rec при великих порогах ($K\ge 50$): HR@100 у
SASRec досягає 0.6313 проти 0.6151 у GRU4Rec.
Це вказує на те, що трансформерна увага точніше охоплює довгострокові
контекстні зв'язки, що корисно для завдань з широким списком рекомендацій.
Проте значно більший час навчання обмежує їх застосовність у
ресурсообмежених середовищах.

\paragraph{NARM.}
NARM (NDCG@10\,=\,0.1163) займає третє місце серед послідовних нейронних
моделей. Механізм уваги дозволяє йому виділяти ключові точки сесії, проте
двоетапна архітектура (GRU\,+\,увага) поступається чистій трансформерній
увазі SASRec за кожним показником.

\paragraph{FPMC як гібридна модель.}
FPMC суттєво перевершує загальні базові моделі, однак значно відстає від
повністю нейронних послідовних методів: NDCG@10\,=\,0.0813 проти 0.1163
у NARM.
Це вказує на обмеженість першопорядкового марківського наближення порівняно
з нелінійними нейронними архітектурами.

% ============================================================
% SECTION 6 — ВИСНОВКИ
% ============================================================
\section{Висновки}

\begin{enumerate}
    \item Послідовна колаборативна фільтрація, що враховує хронологічний
          порядок взаємодій, кардинально переважає статичні методи у
          задачі next-item prediction: середній приріст NDCG@10 становить
          $\approx$4-кратний. Це робить моделювання часової динаміки
          практично обов'язковим для продакшн-рекомендаційних систем.

    \item GRU4Rec є найкращою моделлю за інтегральним критерієм
          «якість\,/\,час навчання»: найвищий NDCG@10\,=\,0.1436 при
          помірних 644~с.
          SASRec та BERT4Rec доцільні тоді, коли пріоритетом є охоплення
          широкого списку ($K\ge 50$), а обчислювальні ресурси дозволяють
          витратити 2\,000--2\,500~с на навчання.

    \item FPMC підтверджує важливість послідовності навіть у спрощеному
          марківському наближенні, однак для досягнення максимальної
          якості необхідні нелінійні нейронні архітектури.
\end{enumerate}

% ============================================================
% REFERENCES
% ============================================================
\newpage
\begin{thebibliography}{10}

\bibitem{harper2015movielens}
F.\,M.~Harper and J.\,A.~Konstan,
``The MovieLens datasets: History and context,''
\textit{ACM Trans. Interact. Intell. Syst.},
vol.~5, no.~4, Art. no.~19, Dec. 2015,
doi:~10.1145/2827872.

\bibitem{zhao2021recbole}
W.\,X.~Zhao et~al.,
``RecBole: Towards a unified, comprehensive and efficient framework
for recommendation algorithms,''
in \textit{Proc. 30th ACM Int. Conf. Inf. Knowl. Manage. (CIKM)},
Virtual Event, QLD, Australia, Nov. 2021, pp.~4653--4664,
doi:~10.1145/3459637.3482016.

\bibitem{akiba2019optuna}
T.~Akiba, S.~Sano, T.~Yanase, T.~Ohta, and M.~Koyama,
``Optuna: A next-generation hyperparameter optimization framework,''
in \textit{Proc. 25th ACM SIGKDD Int. Conf. Knowl. Discovery Data Mining (KDD)},
Anchorage, AK, USA, Aug. 2019, pp.~2623--2631.

\bibitem{rendle2009bpr}
S.~Rendle, C.~Freudenthaler, Z.~Gantner, and L.~Schmidt-Thieme,
``BPR: Bayesian personalized ranking from implicit feedback,''
in \textit{Proc. 25th Conf. Uncertainty Artif. Intell. (UAI)},
Montreal, QC, Canada, Jun. 2009, pp.~452--461.

\bibitem{sarwar2001itemknn}
B.~Sarwar, G.~Karypis, J.~Konstan, and J.~Riedl,
``Item-based collaborative filtering recommendation algorithms,''
in \textit{Proc. 10th Int. Conf. World Wide Web (WWW)},
Hong Kong, May 2001, pp.~285--295,
doi:~10.1145/371920.372071.

\bibitem{rendle2010fpmc}
S.~Rendle, C.~Freudenthaler, and L.~Schmidt-Thieme,
``Factorizing personalized Markov chains for next-basket recommendation,''
in \textit{Proc. 19th Int. Conf. World Wide Web (WWW)},
Apr. 2010, pp.~811--820,
doi:~10.1145/1772690.1772773.

\bibitem{hidasi2016gru4rec}
B.~Hidasi, A.~Karatzoglou, L.~Baltrunas, and D.~Tikk,
``Session-based recommendations with recurrent neural networks,''
in \textit{Proc. 4th Int. Conf. Learn. Represent. (ICLR)},
May 2016.

\bibitem{li2017narm}
J.~Li et~al.,
``Neural attentive session-based recommendation,''
in \textit{Proc. 26th ACM Int. Conf. Inf. Knowl. Manage. (CIKM)},
Singapore, Nov. 2017, pp.~1419--1428,
doi:~10.1145/3132847.3132926.

\bibitem{kang2018sasrec}
W.-C.~Kang and J.~McAuley,
``Self-attentive sequential recommendation,''
in \textit{Proc. IEEE Int. Conf. Data Mining (ICDM)},
Singapore, Nov. 2018, pp.~197--206,
doi:~10.1109/ICDM.2018.00035.

\bibitem{sun2019bert4rec}
F.~Sun et~al.,
``BERT4Rec: Sequential recommendation with bidirectional encoder
representations from transformer,''
in \textit{Proc. 28th ACM Int. Conf. Inf. Knowl. Manage. (CIKM)},
Beijing, China, Nov. 2019, pp.~1441--1450,
doi:~10.1145/3357384.3357895.

\end{thebibliography}

\end{document}
